\subsection{A six-valued obstruction to an intuitionistic logical basis}
\label{sec:ks-obstruction}

An incoming formalism note suggested replacing the four-input
\L ukasiewicz axiom by $K$ and $S$, then attempting to recover classical
logic from the retained quantifier and equality rows.  The following
counterinterpretation resolves that particular logical-bootstrap question
negatively.  It retains all seven other logical rows, both atomic
fresh-variable vacuity schemes, modus ponens, generalization, and
irreflexivity.  It does not include CP or GIC; the conclusion concerns
recovery of the logical basis, not the consistency strength of a complete
modified set-theory search.

\paragraph{Proposition.}
Replace the original single implicational axiom by
\[
 K:\ A\to(B\to A),\qquad
 S:\ (A\to(B\to C))\to((A\to B)\to(A\to C)).
\]
Keep the other seven published logical schemas, including
$\exists_a\forall_a B\to B$, quantified implication distribution,
quantifier commutation, equality existence, and the equality/congruence
rows.  Keep atomic fresh-variable vacuity and IR.
The resulting logical calculus does not derive all instances of Peirce,
double-negation elimination, or the original single axiom.  It also does
not derive unrestricted fresh-variable vacuity.

\paragraph{The algebra.}
Let $H=\{0,1,2\}\times\{0,1\}$ be the product of two finite Heyting
chains.  The order and meet are coordinatewise, the bottom is $(0,0)$,
and the designated top is $(2,1)$.  On a chain with top $t$, implication
is $x\Rightarrow y=t$ if $x\le y$, and $y$ otherwise.  Define
\[
 Q(x,y)=\bigl((1+y)\,[x=2],\ y\bigr),
\]
where the bracket is 1 when the condition holds and 0 otherwise.
The complete table is
\[
\begin{array}{c|rrrrrr}
h&(0,0)&(0,1)&(1,0)&(1,1)&(2,0)&(2,1)\\\hline
Qh&(0,0)&(0,1)&(0,0)&(0,1)&(1,0)&(2,1).
\end{array}
\]
This operation preserves bottom, top, and finite meets.  Its fixed
points are exactly bottom, $e=(0,1)$, and top.  Directly from the formula,
\begin{equation}
 \neg Q\neg Q(x,y)=\bigl(2[x=2],y\bigr)\le(x,y).
 \tag{12}
 \label{eq:ks-rho6}
\end{equation}

\paragraph{The quantified interpretation.}
Use object domain $D=\{0,1\}$.  Atomic equality has value top exactly
when its assigned objects agree, and bottom otherwise.  Let membership
be irreflexive, with either off-diagonal value chosen from
$\{\bot,e,\top\}$.  For environments $s$, set
\[
 \llbracket\forall_i A\rrbracket_s
 =Q\left(\bigwedge_{d\in D}\llbracket A\rrbracket_{s[i:=d]}\right).
\]
Interpret implication in $H$, falsity by bottom, and the source's
derived existential by $\neg\forall_i\neg$.  Literal code 1,
$R(0,0)$, is bottom by irreflexivity.  Code 0 is top, compatibly with
its numerical alias $0\to0=0$.

This is an algebraic interpretation of the emitted inference rules.
It is not ordinary first-order semantics: substitution and general
fresh-variable vacuity are not being assumed as additional rules.

The validation of the rules is uniform in every formula and environment.
$K$ and $S$ hold because $H$ is Heyting; indeed every intuitionistically
valid implication/falsity scheme holds.  Modus ponens preserves top, and
generalization preserves validity because $Q(\top)=\top$.
Meet preservation of $Q$ gives quantified implication distribution.
It also makes both orders of two distinct quantifiers equal to $Q^2$
of the meet over both assignments, proving commutation.  Coincident
quantifier labels give the same expression directly.

For the retained row $\exists_a\forall_a B\to B$, put
$b=\bigwedge_d\llbracket B\rrbracket_{s[a:=d]}$.
The inner universal has value $Qb$, independent of the outer a, so the
existential universal has value $\neg Q\neg Qb\le b\le
\llbracket B\rrbracket_s$ by~\eqref{eq:ks-rho6}.
Equality existence holds because one assignment makes equality top:
the meet of its negations is bottom and $\neg Q\bot=\top$.
Equality congruence and the Euclidean equality row hold by actual
identity of objects, including allowed label collisions.
Fresh-variable quantification leaves an atomic value unchanged before
applying $Q$, and all permitted atomic values are fixed points; hence
both atomic vacuity schemes hold.  IR has value top.

\paragraph{The failure.}
Choose both off-diagonal membership values to be $e$, and distinct
labels $a,b,c$ with an environment assigning different objects to a and b.
Let
\[
 A=R(a,b),\qquad P=\forall_c\neg A.
\]
Their values are
\[
 \llbracket A\rrbracket=(0,1),\quad
 \llbracket\neg A\rrbracket=(2,0),\quad
 \llbracket P\rrbracket=(1,0).
\]
Consequently
\[
 \llbracket\neg\neg P\to P\rrbracket
 =\llbracket((P\to F)\to P)\to P\rrbracket
 =(1,1)\ne\top.
\]
The same value occurs for $\neg A\to\forall_c\neg A$, although c
is absent from A.  Thus an attempted induction from atomic vacuity
fails already at negation.  The original single axiom also fails with
formula arguments $a=F$, $b=\neg A$, $c=P$, and $u=F$; each argument
is an actual formula in the source language.  Induction on derivations
now proves all claimed nonderivability statements.

\paragraph{Executable certificate and scope.}
The standalone script \texttt{check\_ks\_obstruction.py} checks
11,690 declared cases: the algebra, the scalar identity, quantified
distribution and commutation on arbitrary two-point value vectors,
all equality/membership rows with variable-label collisions, and the
explicit counterexamples.  It runs using only the Python standard
library; explicit exceptions keep every check active under
\texttt{python -O}.  This is an ordinary mathematical proof accompanied
by a finite certificate, not a Lean or Rocq proof.

The result is new relative to the examined incoming notes, without a
claim of novelty relative to the algebraic-logic literature.  CP and
GIC might supply additional principles; their ability to repair this
changed calculus would require a separate, theory-specific argument.
In particular, no nonhalting theorem or smaller binary-machine bound
is inferred from this counterinterpretation.
